% ====================================================================
%  RESPONSE TO REVIEWER 2
%  Manuscript TP-1083 -- "The Incommensurability Principle in
%  Biological Transport" -- Transport Phenomena (De Gruyter Brill)
%  Standalone document; compile with: pdflatex Response-to-Reviewer-2
% ====================================================================
\documentclass[11pt]{article}
\usepackage[a4paper,margin=1in]{geometry}
\usepackage{amsmath,amssymb}
\usepackage[dvipsnames]{xcolor}
\usepackage{enumitem}
\usepackage{hyperref}
\hypersetup{colorlinks=true,linkcolor=red!70!black,urlcolor=red!70!black}

% Reviewer 2 changes are shown in RED in the highlighted manuscript.
\newcommand{\rtwo}[1]{\textcolor{red}{#1}}
\newenvironment{comment}%
  {\medskip\noindent\itshape\color{black!75}\rule{\linewidth}{0.4pt}\par\nobreak Reviewer comment.\par\nobreak}%
  {\par\nobreak\rule{\linewidth}{0.4pt}\medskip}
\newcommand{\response}{\par\medskip\noindent\textbf{Response.}\ }

\begin{document}

\begin{center}
{\large\bfseries Response to Reviewer 2}\\[4pt]
Manuscript \textbf{TP-1083} --- \emph{The Incommensurability Principle in
Biological Transport}\\
\emph{Transport Phenomena} (De Gruyter Brill)
\end{center}

\bigskip
\noindent We thank Reviewer~2 for a detailed and constructive review. The four
points raised have led to substantive additions that make the theory's
assumptions explicit, broaden and sharpen its empirical test, and improve its
accessibility. Below we reproduce each comment (paraphrased from the report) and
detail our response. \textbf{Changes prompted by Reviewer~2 are typeset in
\rtwo{red} in the highlighted copy of the manuscript.} For reference, the
highlighted copy uses a distinct colour per source, as the editor requested: blue
for Reviewer~1, \rtwo{red} for Reviewer~2, green for editorial requests, and
magenta for author-initiated refinements.

\bigskip
\hrule
\bigskip

% --------------------------------------------------------------------
\noindent\textbf{Comment 1 --- Assumptions and alternatives.}
\begin{comment}
The derivations rest on assumptions---evolutionary optimisation, informational
constraints, metabolic fitness landscapes---that are not uniquely guaranteed.
The paper should add an explicit discussion of these assumptions and of the
possible alternatives, even where empirical tests are not always feasible.
\end{comment}

\response We have added a dedicated closing subsection to the Discussion,
``\rtwo{Underlying Assumptions and Alternative Hypotheses}'' (Section~7,
\texttt{sec:assumptions}), that makes explicit the five load-bearing assumptions
and, for each, states the principal alternative and---where one exists---the
observation that would discriminate between them:
\begin{enumerate}[label=(A\arabic*),leftmargin=2.2em,itemsep=2pt,topsep=3pt]
  \item \textbf{Optimality}, vs.\ developmental canalisation, physical
        self-organisation, and neutral evolution. We note that
        self-organisation is not in opposition to optimality---local
        shear-homeostatic remodelling (Hu \& Cai; Ronellenfitsch \& Katifori)
        converges to the same optimum, so it is the \emph{mechanism} and the
        minimax the \emph{target}---whereas the genetic-blueprint alternative is
        the explicit falsification target of the ontogenetic transition test.
  \item \textbf{Minimax vs.\ average-case.} We acknowledge an honest
        non-uniqueness: a broad-prior Bayesian (expected-cost) optimum converges
        to nearly the same $\alpha^*$, so the two are hard to separate with
        present data. We adopt the minimax as the more conservative,
        assumption-light choice, requiring no commitment to a particular
        activity prior.
  \item \textbf{Informational bound}, vs.\ long-range biological coordination or
        genetic pre-specification---the former constrained by the phase-lag
        blind spot of Section~2.
  \item \textbf{Metabolic cost currency}, vs.\ other selective currencies
        (developmental robustness, fatigue life, thermoregulation, repair,
        reproduction). This is mitigated by the parameter-independence
        (Topological Rigidity) result, and multi-currency organs (e.g.\ the
        kidney) are shown to fall \emph{outside} the two-channel minimax by
        construction.
  \item \textbf{Incommensurability itself}, vs.\ a single unified energy budget;
        its empirical signature is the dual-attractor structure.
\end{enumerate}
A closing ``Epistemic status'' paragraph states plainly which assumptions are
falsifiable, which are supported, which delimit the domain, and which (A2)
cannot currently be separated from its leading alternative.
\emph{Location:} Section~7, \rtwo{red}.

% --------------------------------------------------------------------
\bigskip
\noindent\textbf{Comment 2 --- Broader empirical validation.}
\begin{comment}
Strengthen and broaden the empirical tests of the Incommensurability Principle
and of the minimax attractor across a wider diversity of biological systems
(beyond the present set of networks).
\end{comment}

\response Rather than padding the validation table with low-resolution or
speculative entries---which would weaken rather than strengthen the case---we
have reframed the existing $27$-network catalogue as what it actually is: a
discriminating, parameter-free, cross-kingdom test of the principle's central
prediction, namely that the realised exponent must track \emph{which}
incommensurable cost channels are physically active. We added (Section~6,
``Independent Datasets'', \rtwo{red}, with a twin framing paragraph in
Supplemental Section~S6) an explicit statement of the sorting rule:
\begin{itemize}[leftmargin=1.4em,itemsep=2pt,topsep=2pt]
  \item non-pulsatile networks---botanical xylem, airways, leaf venation, slime
        mold, mycelium, quasi-static venous return (ten systems, all observed in
        $2.80$--$3.15$)---at the viscous attractor $\alpha_t\to 3.0$;
  \item pulsatile mammalian arteries at the minimax $\alpha^*\approx 2.7$;
  \item wave-dominated or planar ($d=2$) networks near $\alpha_w\to 2.0$.
\end{itemize}
The sorting holds across three biological kingdoms (animal, plant,
fungal/protist), both embedding dimensions, and pathological states---a span no
single-attractor law (Murray's $\alpha=3$ or area-preserving $\alpha=2$ alone)
can reproduce---and is falsifiable in two directions: a non-pulsatile network at
$\alpha\approx 2.7$, or a high-Womersley pulsatile conduit at $\alpha\approx
3.0$, would refute it. We also identify the most informative directions for
extension \emph{beyond} the present set: taxa with divergent cardiac duty cycles
(avian and reptilian circulations), invertebrate open circulatory systems, and
engineered organ-on-chip networks with programmable pulsatility.

% --------------------------------------------------------------------
\bigskip
\noindent\textbf{Comment 3 --- Uniqueness of the linear cost functional
(central point).}
\begin{comment}
The claim that the linear fractional excess is the \emph{unique} admissible cost
functional must be supported more rigorously: explicitly consider alternative
cost formulations, assess their limitations, and justify excluding them in
favour of the linear form.
\end{comment}

\response This is the central methodological point, and we have addressed it
directly. Arguments that were previously scattered through the text (the
exclusion of the logarithm, the weighted-Jensen consistency condition, and
non-substitutability) are now consolidated into a single new Supplemental
subsection, ``\rtwo{Systematic Exclusion of Alternative Cost Functionals}''
(Section~S7, \texttt{sec:S7\_alternatives}). It states four admissibility
requirements,
\begin{enumerate}[label=(\roman*),leftmargin=2.2em,itemsep=1pt,topsep=2pt]
  \item scale invariance,
  \item compositional (weighted-Jensen) consistency,
  \item Onsager linear response near the optimum,
  \item non-substitutability of the two cost channels,
\end{enumerate}
and tests six candidate functionals against them in a summary table
(\texttt{tab:cost\_alternatives}): the linear $x-1$, quadratic $(x-1)^2$,
logarithmic $\ln x$, power-law $x^q-1$, an absolute-threshold form, and a
multiplicative $\mathcal{C}_{\mathrm{wave}}\cdot\mathcal{C}_{\mathrm{transport}}$.
Only the linear excess satisfies all four; each alternative fails at least one
(the quadratic, logarithmic, and power-law forms violate compositional
consistency and Onsager linearity; the absolute-threshold form breaks scale
invariance; the multiplicative form violates non-substitutability). We state
explicitly that linearity is therefore \emph{forced} by the four requirements,
and can be relaxed only by rejecting one of them---which we frame as a concrete
falsification target. A pointer to Section~S7 was added at the
non-substitutability remark in Section~3. \emph{Location:} Section~S7 and
Section~3, \rtwo{red}.

% --------------------------------------------------------------------
\bigskip
\noindent\textbf{Comment 4 --- Expository density.}
\begin{comment}
Several formal constructs in Sections~2--6 are introduced too quickly; expand
the intermediate steps and add deeper biological interpretation of the formal
results.
\end{comment}

\response We have carried out a systematic accessibility pass, working through
the manuscript one section at a time from the beginning and inserting
plain-language ``bridge'' passages wherever a formal construct was previously
introduced too quickly. Each bridge either supplies a missing intermediate step
or gives the biological/physical meaning of a formal result, and adds text
without removing or altering the original derivations. The new passages are:
\begin{itemize}[leftmargin=1.4em,itemsep=3pt,topsep=2pt]
  \item \textbf{Section~1 (Introduction).} A word-level reading of the Kinematic
        Matching theorem (the Womersley number as a ``switch'': viscous resistors
        below $\mathrm{Wo}_c$, pulsatile waveguides above it; two thresholds
        $\Rightarrow$ incommensurability), placed between the theorem statement
        and its proof; and a short gloss of what an \emph{evanescent mode} is.
  \item \textbf{Section~2 (the no-go proposition).} A bricklayer analogy stating,
        in plain terms, what the proposition claims---that a purely local rule
        cannot keep the whole tree scale-free because it would need the
        non-local generation index---which also motivates Step~3 of the proof.
        (The intermediate derivation of the cost ratio $\Lambda(r)\propto r^{2}$
        from the gradient scalings of Step~1 was likewise made explicit.)
  \item \textbf{Section~3 (gauge / ATP).} Three bridges: (i)~reconciling the
        ``absolute'' and ``fractional'' cost statements via within- versus
        across-organism comparison; (ii)~a plain reading of the Onsager-linearity
        theorem (linear in fractional waste, quadratic in geometric deviation);
        and (iii)~a gloss of what the weighted-Jensen equation actually demands
        (boundary-independent scoring of subsystems).
  \item \textbf{Section~4 (architectural invariance).} A ``game against Nature''
        framing of the minimax (the network picks the geometry, an adversary
        picks the worst operating weight, the saddle is the geometry immune to
        that choice), plus a one-line intuition for the envelope-theorem step.
  \item \textbf{Section~5 (single-mechanism limits).} A paragraph giving the
        physiological meaning of the two limits: the static limit as the regime
        of small and non-pulsatile vessels (Murray's law), the wave limit as
        large conduit arteries (area-preserving impedance matching), and the
        interior minimax as the mid-calibre pulsatile beds that morphometry
        actually samples.
  \item \textbf{Section~6 (architectural transition).} A plain-language account
        of the phase decoupling: because the wave responds to the square root of
        the fluid admittance, the fluid and the wave ``switch on'' at different
        body sizes, so they can never match simultaneously---the concrete meaning
        of incommensurability. (The intermediate algebra of the intermediate
        step $\Lambda(r)\propto r^2$ noted above complements this.)
  \item \textbf{Section~9 (retinal paradox).} An explicit up-front statement of
        the paradox (two-dimensional diameter scaling but non-planar bifurcation
        angles within one network) and of its resolution by mechanical--fluidic
        decoupling.
\end{itemize}
We deliberately left Sections~7 and~8 unchanged: as synthesis and conclusion they
are already discursive and carry extensive interpretive remarks, so adding
further exposition there would be padding rather than clarification. In the
Supplemental Material, whose detailed derivations are reference material, we added
a single orienting sentence at the head of the densest sections---the admittance
expansion (S1), the structural-residual decomposition (S3), and the thin-wall
limit (S5)---stating what each establishes and why, without altering any
derivation. \emph{Location:} Sections~1--6 and~9 of the main text, and
Supplemental Sections~S1, S3, S5, \rtwo{red}.

% --------------------------------------------------------------------
\bigskip
\hrule
\bigskip
\noindent\textbf{Closing.} These changes leave the manuscript clearer and better
bounded. Beyond the reviewer-driven changes
(\rtwo{red}), the highlighted copy also marks a set of author-initiated
refinements (magenta), including a corrected citation (the
Fåhræus--Lindqvist viscosity figure, re-stated to match the source) and
consistent disambiguation of previously overloaded symbols ($\eta_d$, the
cardiac duty cycle, vs.\ $\eta$, the cost-mixing weight; and $\beta$, the
branching ratio, vs.\ $\beta_M$, the metabolic scaling exponent). We have also
tightened the existence proof of the minimax saddle in Section~4: because the
network Lagrangian is convex in the geometric exponent and linear in the
cost-mixing weight, the existence and uniqueness of the saddle now follow from an
elementary convex-envelope argument, and the previously cited general
(quasiconvex--quasiconcave) minimax theorem is no longer required.

\end{document}
